% PNP RESEARCH REPORT::CAPPED_BUILD_BEACON::2026-10-07
\documentclass[11pt]{article}
\usepackage[margin=1in]{geometry}
\usepackage{amsmath,amssymb,hyperref}
\title{Executable Capped WRAPPER Portfolio and Recovered Random-Bitstream Controls}
\date{7 October 2026}
\begin{document}\maketitle
\paragraph{Result.} An executable sequential capped-restart portfolio was built
around the existing double-wrapper implementations. It passes the bounded
correctness and failure-handling tests. It does not provide the joint resource
bound needed for the P-versus-NP argument. P versus NP remains OPEN.

\paragraph{Construction.} The standalone Python runner takes the existing
double-wrapper source as an argument. It runs flat, component and fast-separator
backends in separate Linux processes, one at a time. Every attempt has an enforced
64 MiB process address-space ceiling and a two-second wall timeout. Python
line-event budgets increase geometrically from 64 through 65,536. A terminated
or capped attempt contributes no answer. Only a normally completed exact backend
may return SAT or UNSAT. The finite schedule may return unknown outside the test
scope. No claim is made that this bounded runner decides every SAT input.

All restarted/discarded attempt wall times and reported line events are counted.
Process startup, imports and input parsing enter elapsed time. Address-space
limits include interpreter and process input storage. Parent memory, the OS,
and stored experimental logs are additional resources. Python line events are
an execution budget, not a bit-operation count: C builtins may do substantial
work between events. The wall limit covers such work operationally.

\paragraph{Testing.} Ten cases include empty and contradictory formulas,
independent components, pigeonhole UNSAT, satisfiable and unsatisfiable XOR,
and four seeded random eight-variable formulas. Each of three backends and the
portfolio agrees with an independent truth-table answer, for 40 checks and zero
errors. Production reference-SAT entry points are replaced by raising guards.
A deliberately memory-hungry exact-if-completed wrapper is placed before the
normal candidates. Its allocation fails under the actual OS cap; later solvers
still return correct answers. Fuel aborts also occur, so restart behavior is
exercised. Each batch finishes in under 30 seconds with a 25-second alarm.

An initial instrumentation check found that opcode tracing produced zero events
in this environment. The runner was corrected to line-event fuel and now asserts
positive counts in every direct backend test. The published JSON is from the
corrected runner, with code hashes. No opcode-budget result is claimed.

\paragraph{Why this does not yet instantiate the theorem.}
The existing backends parse and retain full parent CNFs, intermediate formulas
and caches. On families without immediate simplification, a current-scale root
representation remains live. A scheduler can reject excess memory use but cannot
turn that implementation into the implicit-input low-space interface required
by the strongest composition theorem. A fixed experimental 64 MiB cap is not a
proof of $N^{o(1)}$ or sublinear space. On sufficiently large inputs this cap
can instead reject all members. No favorable general-input member is established.
Flat won every tested portfolio race; this run gives no empirical advantage
over selecting flat on this small corpus. The result is scheduler correctness
and enforceable operational caps, not asymptotic improvement.

\paragraph{Recovered random-bitstream construction.}
The primary Search-G precommitment source specifies an IID fair bitstream,
a pre-cut execution record $X$, retained state $P(X)$ and side information
$Z(X)$, followed by a fresh block $B$ selecting a delayed query $f_B(X)$.
Its exact response-signature theorem states
\[
s\ge\log_2\max_z|\{(f_B(X))_B:Z(X)=z\}|,
\]
provided no new $X$-dependent information arrives after the cut. If less is
retained, some reacquisition is necessary. That theorem is already useful and
does not require cryptographic hardness. It does not lower-bound the cost of
reacquisition when the original source remains readable.

Conversation retrieval also located ECPC and PRIME-RATCHET descriptions:
sample fresh stream data at computation milestones, bind step counts and prior
commitments, and freeze provenance at SAT completion. These are discovery leads;
this run did not recover a full primary implementation of those named variants.
The tested controls below are explicitly simplified countermodels, not replacements
for unavailable historical implementations.

\paragraph{Fresh exhaustive controls.}
All 256 histories of two independent four-bit blocks were enumerated. A fixed
guess predicts one of 16 first tokens. After revelation, a passive observer
reproduces all 256 simple transcripts. Exact comparison with a fixed commitment
rejects all 256 wrong commitments. If token identity is instead existentially
free, every case permits an alternative-token extension. Thus fixed identity,
freshness, and existential transcript consistency are different requirements.

A separate Search-G coordinate audit uses 16 four-bit records. Their complete
response signatures are distinct. Retaining only the low two bits creates 24
pairs that some later coordinate challenge distinguishes. Nevertheless all 64
record-coordinate responses are recovered with one source-bit probe when source
reread is allowed. This exactly exhibits necessary reacquisition without costly
reacquisition. No superpolynomial inference follows from either finite control.

\paragraph{Historical hostile results.}
Fresh primary reads of REACQ1 and RANDCLOSE1 recover stronger prior warnings.
REACQ1 charges deterministic same-selector rerun plus audit evaluation and notes
that polynomially many such audits remain polynomial under its premises.
RANDCLOSE1 treats bounded-error polynomial-depth static computations with fresh
bits and adaptive NP calls; under $P=NP$ it places their language in P.
These statements do not prove a small quantitative exponent, so they do not
alone exclude every possible resource-improvement use of randomness. They do
block freshness alone as a superpolynomial contradiction. A real-time beacon
or authenticated timestamp supplies an environmental/protocol resource, not
automatically a lower bound in the standard static SAT computation model.

\paragraph{Sources and scope.}
Repository sources read: the original Search-G precommitment spine, TASK-007,
LSPV-4 REACQ1 (2026-10-01 15:46), RANDCLOSE1 (2026-10-01 14:46), and
WRITE-0003 random provenance audit (2026-10-05 19:44). Previous resource-bound
and double-wrapper source reads remain dependencies. No new external theorem
or nonrelativizing/nonalgebrizing step is asserted. These are computational
tests and source reconciliation; independent review remains pending.

\paragraph{Artifacts.}
\texttt{experiments/proof\_spines/capped\_wrapper\_portfolio.py};
\texttt{results/raw/proof\_spines/2026-10-07\_capped\_portfolio.json};
\texttt{results/raw/proof\_spines/2026-10-07\_beacon\_controls.json}.
The runner's test and beacon modes are separate bounded batches. Reproduce by
passing \texttt{test}, the double-wrapper source path, and an output JSON path;
or \texttt{beacon} and an output path.

\paragraph{Next concrete target.}
Replace full parent retention in one backend by a source-backed input interface,
then charge every nested access and remaining live state. If pursuing Search-G
instead, specify a SAT-required delayed query whose post-revelation cost defeats
source reread and same-selector rerun. Adding further random labels alone does
not discharge that requirement.
% END_OF_VERIFIED_REPORT::LEAD::CAPPED_BUILD_BEACON
\end{document}
